\documentclass{beamer}

\input{MathMacros}

\usepackage{beamerthemesplit}
\usepackage{enumerate}
\usepackage{amssymb}
\usepackage[all,cmtip]{xy}
\usepackage[mathscr]{eucal}
\usepackage[natbib=true,style=authoryear,backend=bibtex,useprefix=true]{biblatex}
\addbibresource{reading.bib}
\usepackage{graphicx,color}

\title{The Icosahedron Experiment: What Actually Works in Practice}
\author{Reuben Brasher}
\date{\today}

\begin{document}

\frame{\titlepage}

\section[Outline]{}
\frame{\tableofcontents}

\section{Lecture 4}

\frame{
\frametitle{Empirical, Not Theoretical}
This lecture is grounded in sustained experimentation rather than theory.
\begin{itemize}
\item Eleven structured working sessions were conducted.
\item Multiple domains were explored including probability and systems.
\item Iterative refinement occurred across attempts.
\item Approaches were continuously compared.
\item Failures were analyzed rather than ignored.
\item Focus remained on what works in practice.
\end{itemize}
}

\frame{
\frametitle{Purpose of the Experiment}
The goal of the experiment was evaluation rather than exploration.
\begin{itemize}
\item Identify repeatable success patterns.
\item Identify consistent failure modes.
\item Distinguish illusion of correctness from real correctness.
\item Move from anecdotal success to systematic method.
\item Replace intuition with evidence.
\item Establish operational principles.
\end{itemize}
}

\frame{
\frametitle{Initial Expectations}
Initial assumptions reflected common intuitions about LLMs.
\begin{itemize}
\item Provide a clear description of the task.
\item Receive working code as output.
\item Expect minimal iteration.
\item Assume internal reasoning by the model.
\item Expect outputs to be reliable.
\item Treat prompting as sufficient.
\end{itemize}
}

\frame{
\frametitle{Observed Reality}
Actual behavior diverged significantly from expectations.
\begin{itemize}
\item Outputs were inconsistent across runs.
\item Logical errors appeared frequently.
\item Systems were brittle under small changes.
\item Outputs appeared correct but were not.
\item Internal reasoning was opaque.
\item Small ambiguities caused large failures.
\end{itemize}
}

\frame{
\frametitle{Core Insight from Early Failures}
Early failures revealed a central principle about system behavior.
\begin{itemize}
\item Lack of constraints leads to variability.
\item Ambiguity propagates through outputs.
\item Models fill gaps with plausible logic.
\item Surface fluency hides deeper errors.
\item Prompting alone cannot ensure reliability.
\item Structure is required for stability.
\end{itemize}
}

\frame{
\frametitle{Failure Case 1: Procedural Logic}
A simple generator exposed weaknesses in unstructured interaction.
\begin{itemize}
\item Probability distributions were incorrect.
\item Validation outputs were fabricated.
\item Surface correctness masked deeper flaws.
\item Reasoning was inconsistent across attempts.
\item Outputs varied unpredictably.
\item Results lacked reproducibility.
\end{itemize}
}

\frame{
\frametitle{Diagnosis of Failure Case 1}
The failure resulted from absence of structure rather than complexity.
\begin{itemize}
\item No explicit schema defined behavior.
\item Domain rules were not grounded.
\item Constraints were not specified.
\item No intermediate validation occurred.
\item Logic was inferred instead of defined.
\item Ambiguity remained unresolved.
\end{itemize}
}

\frame{
\frametitle{Failure Case 2: Hierarchical Complexity}
A more complex system revealed deeper structural issues.
\begin{itemize}
\item Class hierarchies were inconsistent.
\item JSON and code representations diverged.
\item Tests failed due to mismatches.
\item Representations became misaligned.
\item Fragility increased with complexity.
\item System coherence degraded over time.
\end{itemize}
}

\frame{
\frametitle{Diagnosis of Failure Case 2}
The primary issue was misalignment across representations.
\begin{itemize}
\item No unified schema linked layers.
\item Representations evolved independently.
\item Consistency was not maintained.
\item No single source of truth existed.
\item Errors amplified with complexity.
\item Alignment was not enforced.
\end{itemize}
}

\frame{
\frametitle{Transition: What Works Instead}
Failures revealed a consistent pattern of success and failure.
\begin{itemize}
\item Unstructured generation produces failure.
\item Structured transformation produces reliability.
\item Variability decreases with constraint.
\item Explicit steps improve control.
\item Structure enables validation.
\item The approach shifted to scaffolding.
\end{itemize}
}

\frame{
\frametitle{Success Case: Layered Scaffolding}
A structured workflow demonstrated stable performance.
\begin{itemize}
\item Natural language defined the domain.
\item Concepts were expressed as tables.
\item Tables were converted to JSON.
\item JSON was mapped to typed models.
\item Models produced implementation code.
\item Code was validated with tests.
\end{itemize}
}

\frame{
\frametitle{Key Insight from Success}
Successful systems shared a common structural property.
\begin{itemize}
\item Each layer constrained the next.
\item Ambiguity decreased step by step.
\item Representations remained aligned.
\item Errors were detected early.
\item Outputs became predictable.
\item Stability emerged from structure.
\end{itemize}
}

\frame{
\frametitle{The Scaffolding Model}
The experiment formalized a repeatable method.
\begin{itemize}
\item Priming defines domain and vocabulary.
\item Grounding introduces structured representations.
\item Decomposition breaks tasks into parts.
\item Each stage constrains the next.
\item Steps are sequenced explicitly.
\item The process becomes repeatable.
\end{itemize}
}

\frame{
\frametitle{Scaffolding Is Required}
Scaffolding is necessary rather than optional.
\begin{itemize}
\item Without priming the domain is misunderstood.
\item Without grounding ambiguity persists.
\item Without decomposition complexity overwhelms.
\item All components are required.
\item Skipping steps reintroduces failure.
\item Reliability depends on structure.
\end{itemize}
}

\frame{
\frametitle{LLM as Pair Programmer}
The model behaves similarly to a junior developer.
\begin{itemize}
\item It is capable but inconsistent.
\item It requires supervision.
\item It performs well with clear instructions.
\item It struggles with ambiguity.
\item It benefits from iterative feedback.
\item It does not ensure correctness alone.
\end{itemize}
}

\frame{
\frametitle{Role Mapping}
Effective collaboration follows structured roles.
\begin{itemize}
\item The human defines direction.
\item The model executes transformations.
\item Roles may shift depending on context.
\item Feedback is continuous.
\item Validation remains human responsibility.
\item Oversight is always required.
\end{itemize}
}

\frame{
\frametitle{Practical Workflow}
A consistent workflow emerged from repeated sessions.
\begin{itemize}
\item Define the schema explicitly.
\item Generate structured representations.
\item Produce implementation code.
\item Generate validation tests.
\item Review outputs carefully.
\item Refine iteratively.
\end{itemize}
}

\frame{
\frametitle{Critical Constraint}
Large tasks must be decomposed into smaller steps.
\begin{itemize}
\item Complexity overwhelms the model.
\item Errors compound without visibility.
\item Validation becomes impossible.
\item Decomposition enables control.
\item Smaller steps improve accuracy.
\item Structure enables correctness.
\end{itemize}
}

\frame{
\frametitle{Connection to System Design}
The findings align with established engineering principles.
\begin{itemize}
\item Systems require explicit structure.
\item Boundaries reduce complexity.
\item Representations must align.
\item Determinism improves reliability.
\item Separation of concerns is essential.
\item Architecture governs behavior.
\end{itemize}
}

\frame{
\frametitle{Implications for Developers}
Benefits vary depending on experience level.
\begin{itemize}
\item Experienced developers gain productivity.
\item They can evaluate outputs effectively.
\item They understand decomposition.
\item They recognize subtle errors.
\item They enforce constraints.
\item They control system behavior.
\end{itemize}
}

\frame{
\frametitle{Risks for Novices}
Inexperience increases the likelihood of failure.
\begin{itemize}
\item Over-trust in generated outputs occurs.
\item Logical errors go undetected.
\item Workflows remain unstructured.
\item Plausible results are misinterpreted.
\item Systems become fragile.
\item Reliability decreases significantly.
\end{itemize}
}

\frame{
\frametitle{Required Skills}
Effective use requires new engineering competencies.
\begin{itemize}
\item Prompt structuring is necessary.
\item Representation design is critical.
\item Task decomposition is required.
\item Output evaluation is essential.
\item Iterative refinement is standard.
\item Skills focus on engineering discipline.
\end{itemize}
}

\frame{
\frametitle{Synthesis of the Experiment}
The experiment leads to a clear conclusion.
\begin{itemize}
\item LLMs are not reliable by default.
\item Reliability comes from structured workflows.
\item Scaffolding stabilizes behavior.
\item Engineering replaces prompting.
\item Systems outperform ad hoc interaction.
\item Structure determines success.
\end{itemize}
}

\frame{
\frametitle{Transition to Phase 3}
Even structured systems encounter limitations at scale.
\begin{itemize}
\item Failures still occur under complexity.
\item Scaling introduces new risks.
\item Reliability becomes systemic.
\item Errors become harder to detect.
\item Systems evolve over time.
\item Next focus is failure at scale.
\end{itemize}
}

\begin{frame}[t,allowframebreaks]
\frametitle{References}
\printbibliography
\end{frame}

\end{document}